\begin{afterword}
\summaryhead

The post asks why some people are troubled by the idea that not every possible mind would agree with us. Mind design space is vast, so for any argument there would seem to be some possible mind it does not persuade. The hope for universal arguments, the author says, comes from picturing a ghost in the machine: a core of reasonableness above the code that any valid argument would reach. But the code is the AI. Physically, for any system that assents to an argument one can build a similar system that does not, so compulsion is a property of minds, not of arguments.

The same answers the charge that Bayesian priors are arbitrary. The ideal philosopher who could be persuaded from no assumptions is the ghost again; a mind with every assumption removed is a rock. So validity, rationality and objectivity cannot rest on universal compulsion. The post ends with two errors about Friendly AI: expecting any AI to reach the one true morality on its own, and hoping a perfectly free AI will reach virtue.

\responsehead

The post's central claim, as it states it, is correct: for any argument some possible mind is not moved by it, and ``compulsion is not a property of arguments.'' The best support is the constructive case of the transistor and the little grey man, which needs no count of minds. For the design of AI the point has practical force, and Nick Bostrom later stated the related claim, that intelligence and final goals can vary independently, as the orthogonality thesis.

The weakness is in the views the post answers. It describes its opponents as people for whom disagreement that persists in principle would make a claim ``arbitrary,'' and it names no philosopher who says so. The classic view that some laws bind every rational mind does not claim that every mind is in fact moved. Kant held that a moral law is valid for all rational beings, and called it an ``ought'' because the human will is not necessarily determined by it. On a view like his a mind that is not moved is irrational, not a counterexample, and the post's conclusion, that validity ``cannot rely on'' universal compulsion, is consistent with it. Whether anything binds all rational minds in that sense is the question that matters, and the post leaves it for later.

The same shift appears with Bayesian priors. The post says ``many philosophers'' call Bayesianism arbitrary because no argument can force every possible journal editor to agree. As the Stanford Encyclopedia describes it, the common worry is that subjective Bayesianism permits too many priors, including ones that defy induction. That is a question about which priors are reasonable. The post raises it itself when it calls one prior perhaps less reasonable, and the answer about ghosts and rocks does not address it.

The people the post does name are not quoted. Two are said to ``dream of commands that no sufficiently advanced mind can disobey,'' a motive for which the post gives no evidence, and the last paragraph offers John C. Wright's conversion as a warning without saying how it bears on his argument. The post does include the author's own younger self among those who made the milder error.

In short: a correct and well-illustrated point about possible minds, aimed at opponents who go unnamed or unquoted, while the classic view it bears on, Kant's, does not claim that every rational being is in fact moved by the moral law.
\end{afterword}
